\section{INTRODUCTION}
\label{sec:introduction}

\subsection{Background}

The Portable Document Format is the dominant container for
presentation-critical text---textbooks, scientific articles, contracts and
technical reports---precisely because it fixes the visual result: glyph
positions, fonts, tables, mathematical expressions and figures are pinned to a
coordinate system that renders identically on every device~\citep{pdf-spec}.
Over the past decade, considerable research effort has been directed at the two
capabilities that a document translation system must combine. On the one hand,
document layout analysis has matured from heuristic segmentation into a
learned detection problem, supported by large human-annotated corpora such as
DocLayNet~\citep{doclaynet} and by detectors designed for the extreme scale
variation of document pages~\citep{doclayout-yolo}; surveys of the field
consistently report that the quality of this stage propagates directly into
every downstream task~\citep{doclayout-survey}. On the other hand, large
language models (LLMs) have substantially changed machine translation,
particularly for lower-resource pairs and for text that needs domain context.
Unlike neural machine translation systems trained for a fixed set of language
pairs~\citep{nmt-survey}, an LLM is multilingual by construction, can exploit a
wide context window, follows instructions, and can be steered towards
consistent terminology through the prompt~\citep{gpt-mt,llm-multilingual-mt}.
Because most LLMs are consumed through hosted APIs, a document translation
system can now be assembled without training any model of its own.

\subsection{Problem Statement}

The very property that makes PDF attractive makes it hostile to translation. A
PDF page is a stream of drawing operators, not a sequence of paragraphs, so
translating one is simultaneously a machine-translation problem and a
\emph{layout reconstruction} problem. Target text almost never has the length
of the source---Vietnamese and German typically expand relative to English,
while Chinese, Japanese and Korean contract---so writing the translation back
into the original boxes produces overflow, collisions with neighbouring
elements, or font shrinkage severe enough to harm readability. Nor is geometry
the only property that must survive: a translated block inherits the visual
role of the block it replaces, so its typeface, size and \emph{colour}---on
whatever background the source page happens to use---are part of the
reconstruction problem rather than cosmetic detail.

It is tempting to treat the input side as solved. Libraries such as PyMuPDF,
\texttt{pdfminer.six} and \texttt{pdfplumber} read a PDF content stream and
return text almost instantly. In practice they only expose the \emph{text
layer} that happens to be stored in the file, rather than analysing what the
page actually displays, and that layer is unreliable in three recurrent ways.
Scanned PDFs are bitmap images and carry no text layer at all, so optical
character recognition (OCR) is unavoidable. \emph{Searchable} scanned PDFs
carry a hidden OCR layer that may disagree with the visible image, so an
extractor silently returns something other than what the reader sees. And even
in born-digital files, embedded fonts with private encodings, missing
\texttt{ToUnicode} maps, formulas stored as isolated drawing operators and
tables stored as scattered text runs all yield extracted strings that are
corrupted or structurally meaningless.

\subsection{Research Gap}

Which stage of a layout-preserving PDF translation system actually binds its
quality? The question is not rhetorical: its answer decides where design and
research effort belong. Existing systems and document-conversion toolkits have
made substantial progress on different parts of the problem---layout-preserving
overlay for scientific PDFs~\citep{pdf2zh}, document-level segmentation and
terminology~\citep{docutranslate}, OCR-and-rebuild workflows for scanned
input~\citep{retain-pdf}, and PDF-to-structured-output parsers for downstream
LLM use~\citep{docling,marker,mineru,olmocr}. What remains hard to infer from
this literature is not whether these components are useful, but how much each
stage contributes to the failure modes of the final translated document. The
available evidence is difficult to compare directly: some systems export
Markdown or JSON rather than a translated PDF, some prioritise born-digital PDFs
and have limited scanned-document support, and most public evaluations score
either document parsing or plain-text translation rather than the reconstructed
artifact.

Three observations motivate measuring the stages apart. First, \textbf{structure
analysis is lossy in a way translation cannot repair}: a region the parser never
emits is content no downstream model can restore, and a mis-read glyph becomes
the translator's input, so errors here are not additive with translation errors
but place a \emph{ceiling} on the pipeline---yet the stage is almost never scored
against ground truth \emph{as a stage}, because a segmentation error and a
translation error arrive together in the final PDF. Second, \textbf{LLM
translation may be a smaller controllable lever}: a designer chooses
among hosted models but controls neither their training data nor the per-language
resource that dominates their quality, so if competitive models occupy a narrow
band, model selection has little leverage compared with the stages specific to
the document task. Third, \textbf{layout preservation is commonly delegated to
the translator}---instructing the model to keep the target within a length budget
of the source is a widespread device, convenient because it costs nothing at
reconstruction time---yet whether such a budget is honoured, and whether
honouring it relates to quality at all, has not been tested.

\subsection{Hypothesis}
\label{subsec:hypothesis}

This study is organised around a single hypothesis, stated in two parts so that
each part is separately falsifiable.

\paragraph{H1 (the bottleneck).} In layout-preserving PDF translation the
largest controllable headroom may lie in \emph{document structure analysis}
rather than in model choice for translation: the parsing stage can produce
non-recoverable content loss whose incidence is concentrated in identifiable
document slices, whereas competitive LLMs may differ only modestly on the
content that does reach them, and much of the remaining variation may be
explained by the target language rather than by pipeline design.

\paragraph{H2 (where layout must be enforced).} Prompt-level length constraints
are an unreliable proxy for layout fidelity. Even when they influence generation,
they do not tell the system whether rendered text will fit the local geometry of
the page; geometry and appearance must instead be \emph{computed at
reconstruction time from the source page}.

\subsection{The Present Study}

We design and implement a layout-preserving multilingual PDF translation
pipeline whose three stages---document structure analysis, content translation
and document reconstruction---are connected by a single structured intermediate
representation carrying element type, bounding box, page index and reading
order, so that each stage exposes a measurable interface. Two design decisions
follow from the hypothesis. To avoid depending on an embedded text layer that is
absent or unreliable for scanned and searchable-scanned documents, every page is
rendered and analysed visually; the cost of that uniform choice is evaluated at
the parser stage rather than assumed away. To keep layout control in the stage
that can see the page, reconstruction decides both the \emph{size} and the
\emph{colour} of every translated block from the source page itself, rather than
inheriting them from a global default or delegating them to the translator.

The hypothesis is examined through three independent stage-level measurements,
defined in Section~\ref{subsec:method-eval}. A separate final-artifact comparison
is reported in Section~\ref{subsec:res-compliance}.

\begin{description}[leftmargin=0em,labelsep=0.4em,itemsep=0.15em]
  \item[M1] \emph{How much content the parser loses, and where.} The visual
  parsing stage is scored against annotated ground truth on localisation,
  labels, text and reading order, sliced by language, layout and document
  source, with deviations that are \emph{benign} for translation separated from
  those that lose the translator's input---only the latter bound the pipeline
  (H1).
  \item[M2] \emph{How much translation quality depends on model choice.} Five
  LLMs are compared on \texttt{en}$\rightarrow$\texttt{xx} translation over many
  target languages with the pipeline held fixed, contrasting the spread
  \emph{between systems} with the spread \emph{between languages} (H1).
  \item[M3] \emph{Whether a prompt-level length constraint is a useful proxy for
  layout fidelity.} Per-language compliance with the $\pm15\%$ length constraint is
  measured and related to translation quality, contrasting expanding
  (Vietnamese, German) with contracting (CJK) targets (H2).
\end{description}

\subsection{Contributions}

This paper makes three contributions.

\begin{enumerate}[leftmargin=1.4em]
  \item \textbf{A layout-preserving architecture in which reconstruction owns
  geometry and appearance.} Translated text is re-typeset inside the source
  geometry by two decisions computed from the source page rather than from the
  translation: a \emph{size} decision, which derives a canonical font size per
  group of same-role elements from the source text and reduces an individual
  element only when its overflow would actually collide with a neighbour; and a
  \emph{colour} decision, which estimates each element's background from a ring
  of pixels around it and its text colour from the core glyph pixels most
  distant from that background---so the same overlay mechanism serves scanned
  and born-digital pages, and coloured text on non-white backgrounds survives
  translation. Upstream, structurally distinctive content (tables,
  mathematical formulas) is routed through dedicated multimodal verification
  instead of being translated as undifferentiated text
  (Section~\ref{sec:methodology}).

  \item \textbf{A stage-level evaluation protocol that makes parsing
  measurable.} We score the parsing stage directly on positions, labels, text
  and reading order against OmniDocBench~\citep{omnidocbench} instead of
  collapsing its output to Markdown, which would discard exactly the geometric
  information the task depends on, and we report a strict and a
  granularity-invariant matcher \emph{as a pair} so that a deliberate design
  choice is not misread as a detection defect; cell geometry, which OmniDocBench
  cannot support, is scored on PubTables-1M~\citep{pubtables1m}; translation is
  scored on all 55
  language pairs of WMT24++~\citep{wmt24pp} with five LLMs, keeping the
  document-level translation protocol intact
  (Sections~\ref{sec:methodology} and~\ref{sec:results}).

  \item \textbf{Stage-level evidence for H1 and H2.} We show that the dominant
  parser error is benign over-segmentation rather than missed content, but that
  the residual content loss is concentrated and non-recoverable (over half of
  the formulas on three-column pages, one table in five); that the quality
  spread between five LLMs on a fixed language ($\approx 0.05$ COMET-DA) is
  roughly a third of the spread attributable to per-language resource
  ($\approx 0.18$), which script family alone does not explain; and that a
  character-based length constraint is honoured only 69--73\% of the time and is
  essentially uncorrelated with semantic adequacy ($r = 0.066$), i.e.\ it
  measures script density rather than quality (Sections~\ref{sec:results}
  and~\ref{sec:discussion}).
\end{enumerate}

The remainder of the paper is organised as follows.
Section~\ref{sec:literature} reviews the PDF layout problem, the model
families the pipeline builds on, and related translation and document-conversion
systems.
Section~\ref{sec:methodology} presents the proposed architecture and the
experimental design of both evaluations. Section~\ref{sec:results} reports the
results, Section~\ref{sec:discussion} interprets them against H1 and H2 and
states the limitations, and Section~\ref{sec:conclusion} concludes.
